\section{结论与展望}
\label{sec:conclusion}

本文的核心贡献可以概括为一个统一框架、一个架构解法、一组可计算原则、一组范式映射和一项社会技术分析。

\subsection{核心贡献}

\textbf{统一框架：ICA六层模型}
本文提出智能计算体系结构（ICA），将模型原生计算系统定义为六个功能层（L1--L6），
每层有明确的接口契约、设计公理和度量指标。
ICA不是对经典计算机体系结构的机械复制，而是一套面向概率式执行的分层抽象。
它的价值在于将AIOS\cite{mei2024aios}、MemGPT\cite{memgpt2023}、MCP\cite{mcp_intro_2026}、PagedAttention\cite{kwon2023pagedattention}等分散工作统一到同一套分层模型中。

\textbf{架构解法：双平面架构}
本文揭示了"LLM是CPU还是OS"的隐喻冲突根源，提出双平面架构作为统一解答。
概率执行平面（模型推理）对应"能做什么"，确定性控制平面（Agent运行时中的调度器、审批器、日志系统）对应"应该做什么"。
这一分离使得ArbiterOS的"Probabilistic CPU"\cite{arbiteros2025}、AIOS的"kernel"\cite{mei2024aios}和AgentOS的"reasoning kernel"\cite{agentos2025}不再互相矛盾，而是同一系统中不同平面的不同视角。

\textbf{可计算原则：三条量化定律}
语义局部性定律（$S = 1/((1-H)+H\cdot\alpha^{-1})$）、上下文预算定律（$W_{\mathrm{eff}} \leq C \cdot \beta(L)$）和Agent加速比定律（$S_{\mathrm{agent}} = 1/((1-F)+F/(N\cdot E))$）为模型原生计算系统提供了类似Amdahl定律和Roofline模型对传统体系结构那样的可计算原则。
这些定律在形式上与Amdahl定律同构，其价值不在于数学创新，而在于为模型原生计算的设计空间提供了可计算的决策依据。
实证检验表明，这些定律与已发表的实验数据定性一致，但仍需更系统的定量验证。

\textbf{范式映射：从硅到基底的智能计算演进}
本文系统性地将CPU五十年演进的经典教训---Dennard扩展定律与能耗墙、big.LITTLE异构架构与能耗感知调度、专用协处理器分工、ISA抽象与基底无关性---映射到LLM/Agent系统的演进路径。
这一映射不仅提供了历史类比，更提出了两个具有架构意义的原则：\textbf{基底无关性}（智能操作系统应与计算核心的物理实现解耦）和\textbf{木桶原理}（智能系统的可靠性受限于其最弱的能力维度）。

\textbf{社会技术分析：智能计算的产业结构}
本文通过半导体和PC产业的解耦历史，分析了AI栈的结构性动力。权重作为可无限复制的软件这一根本特性，与芯片制造的物理护城河形成鲜明对比，驱动着从垂直整合向fabless AI产业结构收敛的趋势---模型设计商、AI代工厂和OS级平台网关的专业化分工。

\subsection{定位与边界}

本文明确承认以下边界：

第一，这些类比并非严格同构。
模型不是确定性CPU，语义检索不是精确寻址，长期记忆也不是无损分页，Agent运行时也不是成熟操作系统。
但正是这些"不完全相同"，构成了未来研究最值得投入的空间：
如何为语义系统设计新的ABI、如何为不确定推理设计新的OS、如何为长期状态设计新的虚拟内存、如何为多Agent协作设计新的提交协议。

第二，本文的独特贡献不在于第一次把大模型与计算机体系结构进行类比——
Ge等\cite{ge2023llmos}、L2MAC\cite{holt2024l2mac}、Mi等\cite{mi2025llmagentscomputersystems}已经覆盖了几条关键路线。
也不在于第一次讨论KV缓存或记忆管理——
MemGPT\cite{memgpt2023}与MemOS\cite{memos2025}已经把记忆子系统高度系统化。
本文的贡献在于\textbf{第一次将这些互相独立、甚至互相冲突的类比整合为一套全栈模型原生计算架构}，包含统一的分层模型、双平面架构、层间接口契约、设计公理体系和可计算的量化定律。

第三，本文是技术畅想型综述（visionary survey），为未来模型原生计算架构提供研究框架，而非最终方案。
ICA的六层划分和三条定律的参数化形式都可能在实践中被修正和细化。
本文的目的是提出正确的问题和可验证的预测，而非给出终极答案。

第四，本文的框架尚未充分建模人机协作的结构性约束。
工业实证显示，工程师在工作中约60\%使用AI，但只能完全委托0--20\%的任务\cite{anthropic2026agenticreport}。
这一"协作悖论"暗示L5--L6接口的任务提交契约需要增加"委托置信度"维度——反映任务可验证性和组织上下文依赖度——这超出了当前ICA模型的范围。

\subsection{未来方向}

如果说过去十年的重点是"训练出更强的基础模型"，那么未来十年的一个核心问题，可能就是："如何把基础模型、上下文记忆、工具接口与Agent运行时组织成真正可编程、可扩展、可治理的模型原生计算体系结构。"

这个问题的答案可能来自多个方向的交汇：
L2层的推理服务优化（KV缓存、量化、speculative decoding）使模型服务更高效；
L3层的上下文编译和记忆管理使Agent拥有真正的工作记忆；
L4层的协议标准化（MCP、A2A）使工具和Agent的互联从"点对点"进化为"总线式"；
L5层的Agent OS使任务分解、权限控制和失败恢复从"手工编码"进化为"平台原生"；
L6层的工作流自动化使最终用户可以用自然语言定义复杂任务。

经济层面的证据已经令人信服：TELUS创建了超过13{,}000个定制AI解决方案，节省了超过500{,}000小时，工程代码交付速度提升30\%；CRED通过将开发者转移到更高价值工作实现了执行速度翻倍；约27\%的AI辅助工作由原本不会被执行的任务组成\cite{anthropic2026agenticreport}。
这些数据表明，模型原生计算不仅加速现有工作，更拓展了经济可行性的前沿。

在更基本的层面上，第~\ref{sec:paradigm}节的范式讨论论证了智能计算栈必须准备好超越硅基---从神经形态芯片和生物基底到量子处理器---并从纯代码交互扩展到物理世界接口。
这些范式级转变将要求智能ISA契约、评估标准和OS调度模型发生根本性演进。
与此同时，第~\ref{sec:sociotechnical}节的分析表明，这一技术演进正在伴随产业结构的解耦：从垂直整合的AI平台走向专业化分层---fabless模型设计、AI代工厂和OS级平台网关---其模式类似于半导体和PC产业的历史性解耦。

本文为这个方向提供了一个系统化的研究框架：一个分层模型（ICA）、一组设计公理、三条可计算定律、一组范式映射和一个从短期到长期的研究路线图。
希望这一框架能够帮助研究者和工程师在快速演化的模型原生计算生态中，找到系统层面的关键问题和可验证的设计原则。

